\subsection{How $\kappa$ and the controlled pairs are computed, in full}\label{sec:ghow}

The input here is the \emph{replicate-level} table (L2): one row per
$(\text{cell},\text{iteration},\text{method})$, reshaped so that each
$(\text{cell},\text{iteration})$ carries one column per method. Because every
method ran on the same simulated data, a row is a set of paired observations ---
same causal variants, same genotypes, same annotations, same phenotype draw.

\paragraph{Step 1 --- headroom.}
Per row,
\[
h=\mathrm{AP}_{\texttt{polyfun\_oracle}}-\mathrm{AP}_{\texttt{susie}} .
\]
\texttt{susie} is the annotation-blind backbone --- what a method achieves
\emph{without} using annotations at all --- and \texttt{polyfun\_oracle} receives
the simulator's true prior weights, so it is what is achievable \emph{with
perfect} annotation information. Their difference is the room annotations
create on that particular replicate.

\paragraph{Step 2 --- the gate.}
Rows with $h\le\texttt{HEADROOM\_EPS}=0.02$ (or non-finite $h$) are excluded,
because $\kappa$ divides by $h$ and is unstable as $h\to0$. The excluded
fraction and the median headroom are reported rather than dropped silently
(\S\ref{sec:headroom}); on the unannotated arm the gate removes
\emph{everything}, which is correct: the ceiling and the baseline are the same
method there, so $h\equiv0$ and $\kappa$ is undefined by construction.

\paragraph{Step 3 --- $\kappa$.}
On surviving rows, for each focus method $m\in\{$\texttt{functional\_beatrice},
\texttt{fb\_xregion}, \texttt{fb\_pooled}$\}$,
\[
\kappa_m=\frac{\mathrm{AP}_m-\mathrm{AP}_{\texttt{susie}}}{h} .
\]
$\kappa=1$ means $m$ captured all the available headroom; $\kappa=0$ means it
did no better than ignoring the annotations; $\kappa<0$ means it did worse than
the annotation-blind baseline on that replicate.

\paragraph{Step 4 --- the controlled pairs, which are NOT gated.}
Two contrasts are formed on \emph{all} rows, with no headroom gate, because they
are differences rather than ratios and have no denominator to destabilise:
\[
\Delta_{\text{annot}}=\mathrm{AP}_{\texttt{functional\_beatrice}}-\mathrm{AP}_{\texttt{beatrice}},
\qquad
\Delta_{\text{joint}}=\mathrm{AP}_{\texttt{fb\_xregion}}-\mathrm{AP}_{\texttt{functional\_beatrice}} .
\]
Each isolates exactly one mechanism, because the two methods being differenced
are identical in every other respect: same inference network, same
hyperparameters, same data, same seed. $\Delta_{\text{annot}}$ turns the
annotation prior on; $\Delta_{\text{joint}}$ changes that prior from per-region
to shared-across-regions.

\paragraph{Step 5 --- inference on the pairs.}
For a group of $n$ paired rows, $\bar\Delta$, $\mathrm{sd}(\Delta)$,
$\mathrm{SE}=\mathrm{sd}/\sqrt n$ and $t=\bar\Delta/\mathrm{SE}$. The pairing is
what supplies the power: the between-scenario variance that dominates marginal
comparisons cancels exactly, leaving only the within-replicate difference. This
is why $\S\ref{sec:senseg}$ resolves effects of $0.02$ AP decisively while
\S\ref{sec:sensed} cannot resolve a $0.13$ gap between two \emph{different}
methods in a single cell --- $n$ differs by three orders of magnitude and the
nuisance variance has been removed.

\paragraph{Step 6 --- conditional breakdowns.}
The figures in \S\ref{sec:where} are $\bar\Delta$ recomputed within levels of one
design factor at a time, pooling the two annotated strata. $114$ of the $22{,}500$
$\Delta$ values are missing ($0.5\%$), all at $S\in\{5,10\}$, corresponding to
the BEATRICE family's small fit-failure rate; those rows are dropped pairwise.
